
\section{两个引理的证明：第二个}
\label{sec:proof-second}

\emph{\cref{thm:df-drho-estimate}的证明。}这条证明的纲要与\cref{thm:dn-op-estimate}的证明相同：我们证明一条抽象引理，它可以应用于复合 $F_{\rho}=p_{\rho}\circ(\operatorname{Id},\widetilde{\scr{F}}_{\rho})\circ i_{\rho}$。下面就是这条抽象引理。

\begin{lemma}\label{thm:composition-rho-derivative-estimate}
设 $G$、$H$ 与 $K$ 是\trBanachSpace，并设
\[
\varphi_{\rho}:G\longrightarrow H\text{，}\qquad\psi_{\rho}:H\longrightarrow K
\]
是 $\cC^{1}$ 类的（非线性）算子，并且当 $\rho\ge\rho_{0}$ 时对 $\rho$ 可微。我们假设存在 $r_{0}>0$ 与不依赖于 $\rho$ 的常数 $k_{i}>0$（对 $i=1,\dots,7$），使得对每个 $z\in\overline{B}(0,r_{0})\subset G$，都有：
\begin{enumerate}[label=(\arabic*)]
\item \cref{thm:dpsi-dphi-compose} 的诸条件对 $\varphi$ 与 $\psi$ 成立（带有常数 $k_{1},\dots,k_{4}$）。
\item $\Bigl\|\dfrac{\partial\varphi_{\rho}}{\partial\rho}(z)-\dfrac{\partial\varphi_{\rho}}{\partial\rho}(0)\Bigr\|\le k_{5}\|z\|$。
\item $\Bigl\|\dfrac{\partial\psi_{\rho}}{\partial\rho}(\varphi_{\rho}(z))-\dfrac{\partial\psi_{\rho}}{\partial\rho}(\varphi_{\rho}(0))\Bigr\|\le k_{6}\,\|\varphi_{\rho}(z)-\varphi_{\rho}(0)\|$。
\item $\Bigl\|\dfrac{\partial\varphi_{\rho}}{\partial\rho}(0)\Bigr\|\le k_{7}$。\footnote{译注：原书此式分子印作 $\varphi_{\rho}$，疑漏偏导号 $\partial$，此处径改 $\dfrac{\partial\varphi_{\rho}}{\partial\rho}$；上一条（3）的第二项 $\dfrac{\partial\psi_{\rho}}{\partial\rho}(\varphi_{\rho}(0))$ 同类（原书分子亦缺 $\partial$，印作 $\psi_{\rho}$），一并径改。——译者注}
\end{enumerate}
那么存在一个不依赖于 $\rho$ 的 $M>0$，使得对每个 $z\in\overline{B}(0,r_{0})$，我们都有
\[
\Bigl\|\frac{\partial}{\partial\rho}(\psi_{\rho}\circ\varphi_{\rho})(z)-\frac{\partial}{\partial\rho}(\psi_{\rho}\circ\varphi_{\rho})(0)\Bigr\|\le M\|z\|\text{。}
\]
\end{lemma}

\begin{proof}
我们有
\[
\begin{aligned}
&\Bigl\|\frac{\partial}{\partial\rho}(\psi_{\rho}\circ\varphi_{\rho})(z)-\frac{\partial}{\partial\rho}(\psi_{\rho}\circ\varphi_{\rho})(0)\Bigr\|\\
&\qquad\le\Bigl\|\frac{\partial\psi_{\rho}}{\partial\rho}(\varphi_{\rho}(z))-\frac{\partial\psi_{\rho}}{\partial\rho}(\varphi_{\rho}(0))\Bigr\|\\
&\qquad\quad+\Bigl\|(d\psi_{\rho})_{\varphi_{\rho}(z)}\Bigl(\frac{\partial\varphi_{\rho}}{\partial\rho}(z)\Bigr)-(d\psi_{\rho})_{\varphi_{\rho}(0)}\Bigl(\frac{\partial\varphi_{\rho}}{\partial\rho}(z)\Bigr)\Bigr\|\\
&\qquad\le k_{6}\,\|\varphi_{\rho}(z)-\varphi_{\rho}(0)\|\\
&\qquad\quad+\Bigl\|(d\psi_{\rho})_{\varphi_{\rho}(z)}\Bigl(\frac{\partial\varphi_{\rho}}{\partial\rho}(z)\Bigr)-(d\psi_{\rho})_{\varphi_{\rho}(0)}\Bigl(\frac{\partial\varphi_{\rho}}{\partial\rho}(z)\Bigr)\Bigr\|\\
&\qquad\quad+\Bigl\|(d\psi_{\rho})_{\varphi_{\rho}(0)}\Bigl(\frac{\partial\varphi_{\rho}}{\partial\rho}(z)-\frac{\partial\varphi_{\rho}}{\partial\rho}(0)\Bigr)\Bigr\|\\
&\qquad\le k_{6}\,\|\varphi_{\rho}(z)-\varphi_{\rho}(0)\|\\
&\qquad\quad+k_{3}\,\|\varphi_{\rho}(z)-\varphi_{\rho}(0)\|\Bigl\|\frac{\partial\varphi_{\rho}}{\partial\rho}(z)\Bigr\|+k_{4}\Bigl\|\frac{\partial\varphi_{\rho}}{\partial\rho}(z)-\frac{\partial\varphi_{\rho}}{\partial\rho}(0)\Bigr\|\\
&\qquad\le k_{6}\,\|\varphi_{\rho}(z)-\varphi_{\rho}(0)\|\\
&\qquad\quad+k_{3}\,\|\varphi_{\rho}(z)-\varphi_{\rho}(0)\|\left(\Bigl\|\frac{\partial\varphi_{\rho}}{\partial\rho}(z)-\frac{\partial\varphi_{\rho}}{\partial\rho}(0)\Bigr\|+\Bigl\|\frac{\partial\varphi_{\rho}}{\partial\rho}(0)\Bigr\|\right)+k_{4}k_{5}\|z\|\\
&\qquad\le\|\varphi_{\rho}(z)-\varphi_{\rho}(0)\|\left(k_{6}+k_{3}k_{5}\|z\|+k_{7}\right)+k_{4}k_{5}\|z\|\text{。}
\end{aligned}
\]
其次，\footnote{译注：上一链式末行括号内末项原书印作 $k_{7}$；按其上一行的 $k_{3}\|\varphi_{\rho}(z)-\varphi_{\rho}(0)\|(\,\cdots\,)$ 以假设 (2)、(4) 代入展开，末项应为 $k_{3}k_{7}$。此处照排；以一个统一常数 $M$ 重新吸收后，引理的最终结论不受影响。——译者注}
\[
\begin{aligned}
\|\varphi_{\rho}(z)-\varphi_{\rho}(0)\|&\le\sup_{t\in[0,1]}\|(d\varphi_{\rho})_{tz}\|\,\|z\|\\
&\le\sup_{t\in[0,1]}\left(\|(d\varphi_{\rho})_{tz}-(d\varphi_{\rho})_{0}\|+\|(d\varphi_{\rho})_{0}\|\right)\|z\|\\
&\le\left(k_{1}\|z\|+k_{2}\right)\|z\|\\
&\le\left(k_{1}r_{0}+k_{2}\right)\|z\|\text{。}
\end{aligned}
\]
这就结束了\cref{thm:composition-rho-derivative-estimate}的证明。
\end{proof}

我们将把它应用于下列情形：
\begin{itemize}
\item $\varphi=i_{\rho}$ 与 $\psi=\operatorname{Id}\times\widetilde{\scr{F}}_{\rho}$；
\item $\varphi=\operatorname{Id}\times\widetilde{\scr{F}}_{\rho}$ 与 $\psi=p_{\rho}$。
\end{itemize}

我们来验证：在两种情形下，这些算子都满足\cref{thm:composition-rho-derivative-estimate}的诸条件：
\begin{itemize}
\item 条件 (1) 已经在\cref{thm:dn-op-estimate}那里得到过验证。
\item 我们来对 $\varphi_{\rho}=i_{\rho}$ 验证条件 (2) 与 (4)。我们有 $i_{\rho}(Y)=\sum y_{i}Z_{i}^{\rho}$，从而
\[
\frac{\partial i_{\rho}}{\partial\rho}(Y)=\sum_{i=1}^{2n}y_{i}\frac{\partial Z_{i}^{\rho}}{\partial\rho}\text{，}\qquad\text{特别地 }\frac{\partial i_{\rho}}{\partial\rho}(0)=0\text{。}
\]
这就给出
\[
\begin{aligned}
\Bigl\|\frac{\partial i_{\rho}}{\partial\rho}(Y)-\frac{\partial i_{\rho}}{\partial\rho}(0)\Bigr\|_{W^{1,p}}&=\Bigl\|\sum_{i=1}^{2n}y_{i}\frac{\partial Z_{i}^{\rho}}{\partial\rho}\Bigr\|_{W^{1,p}}\le\sum_{i=1}^{2n}\Bigl\|y_{i}\frac{\partial Z_{i}^{\rho}}{\partial\rho}\Bigr\|_{W^{1,p}}\\
&\le\sum_{i=1}^{2n}\|y_{i}\|_{W^{1,p}}\Bigl\|\frac{\partial Z_{i}^{\rho}}{\partial\rho}\Bigr\|_{\cC^{1}}
\end{aligned}
\]
（再次由\cref{thm:w1p-product-inequality}）。现在，诸 $Z_{i}^{\rho}(s,t)$ 的选取方式（参见\cref{sec:construction-psi-rho}）保证存在一个不依赖于 $\rho$ 的常数 $C>0$，使得对每个 $i=1,\dots,2n$，
\[
\Bigl\|\frac{\partial Z_{i}^{\rho}}{\partial\rho}\Bigr\|_{\cC^{1}}\le C\text{。}
\]
由此推出
\[
\Bigl\|\frac{\partial i_{\rho}}{\partial\rho}(Y)-\frac{\partial i_{\rho}}{\partial\rho}(0)\Bigr\|_{W^{1,p}}\le C\sum_{i=1}^{2n}\|y_{i}\|_{W^{1,p}}\le 2nC\,\|Y\|_{W^{1,p}}\text{。}
\]
条件 (4) 于是得到满足，因为 $\partial i_{\rho}/\partial\rho(0)=0$。
\item 我们现在来对 $\psi_{\rho}=\operatorname{Id}\times\widetilde{\scr{F}}_{\rho}$ 验证条件 (3)。注意到对 $\varphi_{\rho}=i_{\rho}$，我们有 $\varphi_{\rho}(0)=0$，并且
\[
\|\varphi_{\rho}(z)\|\le\|i_{\rho}\|\,\|z\|\le C_{1}r_{0}\text{。}
\]
因此，只需证明：对每个 $X\in B(0,C_{1}r_{0})$，
\[
\Bigl\|\frac{\partial(\operatorname{Id}\times\widetilde{\scr{F}}_{\rho})}{\partial\rho}(X)-\frac{\partial(\operatorname{Id}\times\widetilde{\scr{F}}_{\rho})}{\partial\rho}(0)\Bigr\|_{L^{p}}\le k_{6}\|X\|_{W^{1,p}}\text{，}
\]
也就是说，证明
\[
\Bigl\|\frac{\partial\widetilde{\scr{F}}_{\rho}}{\partial\rho}(X)-\frac{\partial\widetilde{\scr{F}}_{\rho}}{\partial\rho}(0)\Bigr\|_{L^{p}}\le k_{6}\|X\|_{W^{1,p}}\text{。}
\]
回忆一下
\[
\begin{aligned}
\widetilde{\scr{F}}_{\rho}(X)&=\left(\frac{\partial}{\partial s}+\widetilde{J}\frac{\partial}{\partial t}+\operatorname{grad}H_{t}\right)\left(\widetilde{\exp}_{w_{\rho}}(X)\right)\\
&=D_{1}\widetilde{\exp}_{(w_{\rho},X)}\left(\frac{\partial w_{\rho}}{\partial s}\right)+D_{2}\widetilde{\exp}_{(w_{\rho},X)}\left(\frac{\partial X}{\partial s}\right)\\
&\qquad+\widetilde{J}_{\widetilde{\exp}(w_{\rho},X)}\left(D_{1}\widetilde{\exp}_{(w_{\rho},X)}\left(\frac{\partial w_{\rho}}{\partial t}\right)+D_{2}\widetilde{\exp}_{(w_{\rho},X)}\left(\frac{\partial X}{\partial t}\right)\right)\\
&\qquad+\operatorname{grad}H_{t}\bigl(\widetilde{\exp}(w_{\rho},X)\bigr)\text{。}
\end{aligned}
\]
由此推出\footnote{译注：下两个 $D_{1}\widetilde{J}$ 项的首个变元原书均直接印作 $\dfrac{\partial w_{\rho}}{\partial\rho}$；按本式系把 $\widetilde{J}\bigl(\widetilde{\exp}(w_{\rho},X)\bigr)$ 对 $\rho$ 求导，基点 $\widetilde{\exp}(w_{\rho},X)$ 随 $\rho$ 的变动经链式法则应给出 $D_{1}\widetilde{\exp}_{(w_{\rho},X)}\bigl(\dfrac{\partial w_{\rho}}{\partial\rho}\bigr)$——末项 $(d\operatorname{grad}H_{t})$ 中恰已保留该因子，可资内证——疑系排印脱漏。补入该因子后，此项仍受 $D_{1}\widetilde{\exp}$ 与 $w_{\rho}$ 诸导数的一致\trBounded{}性控制，下文的估计照常成立。两处径改。——译者注}
\[
\begin{aligned}
\frac{\partial}{\partial\rho}\widetilde{\scr{F}}_{\rho}(X)&=D_{1}D_{1}\widetilde{\exp}_{(w_{\rho},X)}\left(\frac{\partial w_{\rho}}{\partial\rho},\frac{\partial w_{\rho}}{\partial s}\right)+D_{1}\widetilde{\exp}_{(w_{\rho},X)}\left(\frac{\partial^{2}w_{\rho}}{\partial\rho\partial s}\right)\\
&\qquad+D_{1}D_{2}\widetilde{\exp}_{(w_{\rho},X)}\left(\frac{\partial w_{\rho}}{\partial\rho},\frac{\partial X}{\partial s}\right)\\
&\qquad+D_{1}\widetilde{J}_{\widetilde{\exp}(w_{\rho},X)}\left(D_{1}\widetilde{\exp}_{(w_{\rho},X)}\left(\frac{\partial w_{\rho}}{\partial\rho}\right),D_{1}\widetilde{\exp}_{(w_{\rho},X)}\left(\frac{\partial w_{\rho}}{\partial t}\right)\right)\\
&\qquad+D_{1}\widetilde{J}_{\widetilde{\exp}(w_{\rho},X)}\left(D_{1}\widetilde{\exp}_{(w_{\rho},X)}\left(\frac{\partial w_{\rho}}{\partial\rho}\right),D_{2}\widetilde{\exp}_{(w_{\rho},X)}\left(\frac{\partial X}{\partial t}\right)\right)\\
&\qquad+\widetilde{J}_{\widetilde{\exp}(w_{\rho},X)}\Bigl[D_{1}D_{1}\widetilde{\exp}_{(w_{\rho},X)}\left(\frac{\partial w_{\rho}}{\partial\rho},\frac{\partial w_{\rho}}{\partial t}\right)\\
&\qquad\qquad+D_{1}\widetilde{\exp}_{(w_{\rho},X)}\left(\frac{\partial^{2}w_{\rho}}{\partial\rho\partial t}\right)+D_{1}D_{2}\widetilde{\exp}_{(w_{\rho},X)}\left(\frac{\partial w_{\rho}}{\partial\rho},\frac{\partial X}{\partial t}\right)\Bigr]\\
&\qquad+\bigl(d\operatorname{grad}H_{t}\bigr)_{\widetilde{\exp}(w_{\rho},X)}\left(D_{1}\widetilde{\exp}_{(w_{\rho},X)}\left(\frac{\partial w_{\rho}}{\partial\rho}\right)\right)\text{。}
\end{aligned}
\]
我们固定一个对子 $(s,t)\in\RR\times S^{1}$，并把\cref{thm:a-b-estimates} 应用于这个公式中出现的所有线性映射 $A_{(p,X)}(\cdot)$ 与双线性映射 $B_{(p,X)}(\cdot,\cdot)$。这里 $p=w_{\rho}(s,t)$ 而 $X=X(s,t)$，并且我们有 $\|X(s,t)\|_{\cC^{0}}\le r_{1}$，这便容许我们使用\cref{thm:a-b-estimates}。在点 $(s,t)$ 处、对 $\RR^{m}$ 的\trEuclidean{}范数，我们得到估计
\[
\begin{aligned}
&\Bigl\|\frac{\partial}{\partial\rho}\widetilde{\scr{F}}_{\rho}(X)-\frac{\partial}{\partial\rho}\widetilde{\scr{F}}_{\rho}(0)\Bigr\|\\
&\qquad\le C\bigl(\|X\|\bigl\|\frac{\partial w_{\rho}}{\partial\rho}\bigr\|\bigl\|\frac{\partial w_{\rho}}{\partial s}\bigr\|+\|X\|\bigl\|\frac{\partial^{2}w_{\rho}}{\partial\rho\partial s}\bigr\|+\bigl\|\frac{\partial w_{\rho}}{\partial\rho}\bigr\|\bigl\|\frac{\partial X}{\partial s}\bigr\|\\
&\quad\qquad+\|X\|\bigl\|\frac{\partial w_{\rho}}{\partial\rho}\bigr\|\bigl\|\frac{\partial w_{\rho}}{\partial t}\bigr\|+\bigl\|\frac{\partial w_{\rho}}{\partial\rho}\bigr\|\bigl\|\frac{\partial X}{\partial t}\bigr\|+\|X\|\bigl\|\frac{\partial^{2}w_{\rho}}{\partial\rho\partial t}\bigr\|+\|X\|\bigl\|\frac{\partial w_{\rho}}{\partial\rho}\bigr\|\bigr)\text{。}
\end{aligned}
\]
由于 $w_{\rho}$ 的所有导数都被一个常数 $k>0$ 一致控制，这就导致
\[
\begin{aligned}
\Bigl\|\frac{\partial}{\partial\rho}\widetilde{\scr{F}}_{\rho}(X)-\frac{\partial}{\partial\rho}\widetilde{\scr{F}}_{\rho}(0)\Bigr\|&\le C\Bigl(2k^{2}\|X\|+3k\,\|X\|+k\Bigl\|\frac{\partial X}{\partial s}\Bigr\|+k\Bigl\|\frac{\partial X}{\partial t}\Bigr\|\Bigr)\\
&\le C_{1}\left(\|X\|+\Bigl\|\frac{\partial X}{\partial s}\Bigr\|+\Bigl\|\frac{\partial X}{\partial t}\Bigr\|\right)\text{。}
\end{aligned}
\]
让 $(s,t)$ 变动并取 $L^{p}$ 范数，我们得到
\[
\begin{aligned}
\Bigl\|\frac{\partial}{\partial\rho}\widetilde{\scr{F}}_{\rho}(X)-\frac{\partial}{\partial\rho}\widetilde{\scr{F}}_{\rho}(0)\Bigr\|_{L^{p}}&\le C_{1}\left(\|X\|_{L^{p}}+\Bigl\|\frac{\partial X}{\partial s}\Bigr\|_{L^{p}}+\Bigl\|\frac{\partial X}{\partial t}\Bigr\|_{L^{p}}\right)\\
&\le 3C_{1}\,\|X\|_{W^{1,p}}\text{。}
\end{aligned}
\]
\item 我们来对 $\varphi_{\rho}=(\operatorname{Id}\times\widetilde{\scr{F}}_{\rho})\circ i_{\rho}$ 验证诸条件。把\cref{thm:composition-rho-derivative-estimate} 应用于上面研究过的情形，便可见条件 (2) 得到满足。按条件 (4)，由于我们有
\[
\varphi_{\rho}(Y)=\bigl(i_{\rho}(Y),\widetilde{\scr{F}}_{\rho}\circ i_{\rho}(Y)\bigr)\text{，}
\]
导数便是
\[
\frac{\partial\varphi_{\rho}}{\partial\rho}(0)=\left(\frac{\partial i_{\rho}}{\partial\rho}(0),\frac{\partial\widetilde{\scr{F}}_{\rho}}{\partial\rho}(0)+(d\widetilde{\scr{F}}_{\rho})_{0}\left(\frac{\partial i_{\rho}}{\partial\rho}(0)\right)\right)=\left(0,\frac{\partial\widetilde{\scr{F}}_{\rho}}{\partial\rho}(0)\right)\text{。}
\]
其次，
\[
\begin{aligned}
\Bigl\|\frac{\partial\widetilde{\scr{F}}_{\rho}}{\partial\rho}(0)\Bigr\|_{L^{p}}&=\Bigl\|\frac{\partial}{\partial\rho}\left(\frac{\partial w_{\rho}}{\partial s}+\widetilde{J}_{w_{\rho}}\frac{\partial w_{\rho}}{\partial t}+\operatorname{grad}H_{t}(w_{\rho})\right)\Bigr\|_{L^{p}}\\
&=\Bigl\|\frac{\partial^{2}w_{\rho}}{\partial\rho\partial s}+D_{1}\widetilde{J}_{w_{\rho}}\left(\frac{\partial w_{\rho}}{\partial\rho},\frac{\partial w_{\rho}}{\partial t}\right)\\
&\qquad+\widetilde{J}_{w_{\rho}}\frac{\partial^{2}w_{\rho}}{\partial\rho\partial t}+\bigl(d\operatorname{grad}H_{t}\bigr)_{w_{\rho}}\left(\frac{\partial w_{\rho}}{\partial\rho}\right)\Bigr\|_{L^{p}}\\
&\le C\bigl(\bigl\|\frac{\partial^{2}w_{\rho}}{\partial\rho\partial s}\bigr\|_{L^{p}}+\bigl\|\bigl\|\frac{\partial w_{\rho}}{\partial\rho}\bigr\|\,\bigl\|\frac{\partial w_{\rho}}{\partial t}\bigr\|\bigr\|_{L^{p}}+\bigl\|\frac{\partial^{2}w_{\rho}}{\partial\rho\partial t}\bigr\|_{L^{p}}+\bigl\|\frac{\partial w_{\rho}}{\partial\rho}\bigr\|_{L^{p}}\bigr)\text{。}
\end{aligned}
\]
现在，$w_{\rho}$ 的所有导数按 $L^{p}$ 范数都是一致\trBounded{}的。这在 $[-1,1]\times S^{1}$ 上成立（因为它们在 $\cC^{\infty}$ 中一致\trBounded）；对于 $|s|\ge 1$，我们利用 $w_{\rho}(s,t)=u(s+\rho,t)$（或 $v(s-\rho,t)$）这一事实以及 $u$ 与 $v$ 的\trExponentialDecay，后者蕴含它们的导数都属于 $L^{p}$。因此我们有
\[
\Bigl\|\frac{\partial\varphi_{\rho}}{\partial\rho}(0)\Bigr\|\le k\text{，}
\]
其中 $k>0$ 不依赖于 $\rho$。
\item 我们来对 $\psi_{\rho}=P_{\rho}$ 验证条件 (3)。我们有
\[
\varphi_{\rho}(0)=\bigl(0,\widetilde{\scr{F}}_{\rho}(0)\bigr)\text{。}
\]
设 $Y_{0}=\widetilde{\scr{F}}_{\rho}(0)$，让我们计算诸 $p_{\rho}^{i}(X,Y)$（对 $i=1,\dots,2n$）的（关于 $\rho$ 的）导数。我们有
\[
\begin{aligned}
\frac{\partial p_{\rho}^{i}}{\partial\rho}(X,Y)&=\frac{\partial P}{\partial\rho}\bigl(w_{\rho},X,Z_{i}^{\rho},Y\bigr)\\
&=\bigl(D_{1}P\bigr)_{(w_{\rho},X,Z_{i}^{\rho},Y)}\left(\frac{\partial w_{\rho}}{\partial\rho}\right)+P\left(w_{\rho},X,\frac{\partial Z_{i}^{\rho}}{\partial\rho},Y\right)\text{。}
\end{aligned}
\]
我们固定一个对子 $(s,t)\in\RR\times S^{1}$，并把\cref{thm:a-b-estimates} 应用于双线性映射 $P_{(p,X)}(\cdot,\cdot)$ 与三线性映射 $D_{1}P_{(p,X,\cdot,\cdot)}(\cdot)$（取 $p=w_{\rho}(s,t)$ 与 $X=X(s,t)$）。这就给出
\[
\begin{aligned}
\Bigl\|\frac{\partial p_{\rho}^{i}}{\partial\rho}(X,Y)-\frac{\partial p_{\rho}^{i}}{\partial\rho}(0,Y_{0})\Bigr\|&\le C\Bigl(\|X\|\,\|Z_{i}^{\rho}\|\,\|Y_{0}\|\Bigl\|\frac{\partial w_{\rho}}{\partial\rho}\Bigr\|+\|Z_{i}^{\rho}\|\Bigl\|\frac{\partial w_{\rho}}{\partial\rho}\Bigr\|\,\|Y-Y_{0}\|\\
&\qquad\qquad+\|X\|\Bigl\|\frac{\partial Z_{i}^{\rho}}{\partial\rho}\Bigr\|\,\|Y_{0}\|+\Bigl\|\frac{\partial Z_{i}^{\rho}}{\partial\rho}\Bigr\|\,\|Y-Y_{0}\|\Bigr)\text{。}
\end{aligned}
\]
由于 $\|Z_{i}^{\rho}\|$、$\|\partial Z_{i}^{\rho}/\partial\rho\|$、$\|\partial w_{\rho}/\partial\rho\|$ 以及 $\|Y_{0}\|=\|\scr{F}_{w_{\rho}}\|$ 都被一个常数 $k$（它不依赖于 $\rho$、$s$ 或 $t$）一致控制，由此推出
\[
\Bigl\|\frac{\partial p_{\rho}^{i}}{\partial\rho}(X,Y)-\frac{\partial p_{\rho}^{i}}{\partial\rho}(0,Y_{0})\Bigr\|\le 2kC\left(\|X\|+\|Y-Y_{0}\|\right)
\]
而取 $L^{p}$ 范数后，
\[
\begin{aligned}
\Bigl\|\frac{\partial p_{\rho}^{i}}{\partial\rho}(X,Y)-\frac{\partial p_{\rho}^{i}}{\partial\rho}(0,Y_{0})\Bigr\|&\le 2kC\left(\|X\|_{L^{p}}+\|Y-Y_{0}\|_{L^{p}}\right)\\
&\le 2kC\left(\|X\|_{W^{1,p}}+\|Y-Y_{0}\|_{L^{p}}\right)\text{。}
\end{aligned}
\]
$p_{\rho}$ 的每个分量 $p_{\rho}^{i}$ 都满足条件 (3)，从而 $p_{\rho}$ 本身也满足条件 (3)。

\end{itemize}

于是，我们已经验证了所有条件，从而可以应用\cref{thm:composition-rho-derivative-estimate} 来结束\cref{thm:df-drho-estimate}的证明。\noqed
